% !TEX root = ../main.tex
\documentclass[../main.tex]{subfiles}
\begin{document}

\chapter{Estado del arte}
\label{chap:estado_del_arte}

\section{Sistemas de visión artificial embebida}

Los sistemas de visión artificial embebida han evolucionado notablemente en la última década. Las soluciones basadas en microcontroladores de altas prestaciones, como los de la familia STM32H7, permiten capturar y procesar imágenes de baja resolución con latencias aceptables. Sin embargo, cuando se requiere el procesamiento de flujos de vídeo en tiempo real a resoluciones \ac{VGA} o superiores, las \ac{FPGA}s se presentan como la alternativa más adecuada gracias a su capacidad de procesar datos en paralelo y al acceso directo a los pines de entrada/salida sin la mediación de un sistema operativo.

Entre los trabajos más representativos en este ámbito destaca el uso de plataformas Zynq de Xilinx, que combinan un procesador ARM con lógica programable y permiten implementar \textit{pipelines} de procesamiento de imagen con aceleración hardware. No obstante, para aplicaciones donde el objetivo principal es la adquisición y transmisión de imagen sin procesamiento local intensivo, una \ac{FPGA} pura como la Artix-7 ofrece menores costes y una complejidad de diseño más controlable.

\section{Interfaces de alta velocidad hacia PC}

La transmisión de flujos de imagen sin comprimir hacia un PC exige interfaces de alta velocidad. Las alternativas más comunes son USB~3.0, GigE Vision y PCIe. En el contexto de plataformas \ac{FPGA} de bajo coste, la familia FT232H de FTDI ocupa un lugar destacado: ofrece un modo \textit{Synchronous FIFO} de 8 bits que permite alcanzar tasas de transferencia de hasta 60~MB/s sobre USB~2.0~\cite{ft232h_datasheet}, con una interfaz síncrona simple de implementar en \ac{VHDL}.

El modo \textit{Synchronous FIFO} del FT232H opera con un reloj de 60~MHz generado por el propio chip. La FPGA actúa como maestro del bus: para transmitir datos (FPGA$\rightarrow$PC) activa una señal y presenta el dato en el bus de datos; para recibir datos (PC$\rightarrow$FPGA) activa otra señal diferente para que el FT232H conduzca el bus y a continuación aserta otra señal para consumir los bytes disponibles.
\section{Sensores de imagen CMOS}

El sensor MT9V111 de ON~Semiconductor es un sensor CMOS VGA (640$\times$480) que entrega datos en formato YCbCr~4:2:2. Dispone de una interfaz de configuración I2C de dos hilos y una interfaz de datos paralela de 8~bits con señales de sincronismo de trama (\texttt{FVALID}), línea (\texttt{LVALID}) y reloj de píxel (\texttt{PIXCLK})~\cite{mt9v111_datasheet}. El protocolo de lectura de registros del MT9V111 utiliza dos transacciones I2C separadas por una condición de \textit{STOP}: primero una escritura de la dirección del registro y a continuación una lectura del valor, a diferencia del \textit{Repeated START} habitual en otros dispositivos I2C.

\section{Metodologías de verificación de hardware digital}

La verificación de diseños VHDL ha evolucionado desde los bancos de pruebas ad hoc hasta metodologías estructuradas basadas en \textit{frameworks} como UVVM (\textit{Universal VHDL Verification Methodology})~\cite{uvvm}. UVVM proporciona una biblioteca de utilidades (\texttt{uvvm\_util}) con funciones de logging, gestión de alertas y comprobación de valores, así como un framework de \textit{Verification Components} (VVCs) que permite abstraer los agentes de verificación y coordinarlos desde un secuenciador central. La metodología promueve la reutilización de componentes de verificación y facilita la generación automática de informes de cobertura y trazabilidad de requisitos.

\end{document}
